
\documentclass{article}
\usepackage{colortbl}
\usepackage{makecell}
\usepackage{multirow}
\usepackage{supertabular}

\begin{document}

\newcounter{utterance}

\centering \large Interaction Transcript for game `cladder', experiment `full\_v1.5\_default', episode 2060 with llama{-}3.1{-}8b{-}instruct.
\vspace{24pt}

{ \footnotesize  \setcounter{utterance}{1}
\setlength{\tabcolsep}{0pt}
\begin{supertabular}{c@{$\;$}|p{.15\linewidth}@{}p{.15\linewidth}p{.15\linewidth}p{.15\linewidth}p{.15\linewidth}p{.15\linewidth}}
    \# & \multicolumn{2}{c}{Player} && \multicolumn{2}{c}{Game Master} \\
    \hline

    \theutterance \stepcounter{utterance}  
    & & & \multicolumn{4}{p{0.6\linewidth}}{
        \cellcolor[rgb]{0.9,0.9,0.9}{
            \makecell[{{p{\linewidth}}}]{
                \texttt{\tiny{[P1$\langle$GM]}}
                \texttt{You are an expert in causal inference. The following question is not a typical commonsense query, but rather a meticulously designed question created by a professor specializing in causal inference, intended to assess the students' mastery of the course content.} \\
\texttt{Imagine a self{-}contained, hypothetical world with only the following conditions, and without any unmentioned factors or causal relationships: Encouragement level has a direct effect on studying habit and exam score. Studying habit has a direct effect on exam score. The overall probability of encouragement is 91\%. For students who are not encouraged, the probability of high exam score is 72\%. For students who are encouraged, the probability of high exam score is 52\%. Is high exam score more likely than low exam score overall?} \\
\texttt{Start your answer with 'Yes' or 'No', followed by additional reasoning or evidence to support your explanation.} \\
            }
        }
    }
     \\ \\

    \theutterance \stepcounter{utterance}  
    & \multicolumn{4}{p{0.6\linewidth}}{
        \cellcolor[rgb]{0.561,0.737,0.561}{
            \makecell[{{p{\linewidth}}}]{
                \texttt{\tiny{[P1$\rangle$GM]}}
                \texttt{No.} \\
\\ 
\texttt{To answer this question, we need to apply the principles of causal inference, specifically the concept of counterfactuals and potential outcomes. Let's break down the given information:} \\
\\ 
\texttt{1. Encouragement level has a direct effect on studying habit and exam score.} \\
\texttt{2. Studying habit has a direct effect on exam score.} \\
\texttt{3. The overall probability of encouragement is 91\%.} \\
\texttt{4. For students who are not encouraged, the probability of high exam score is 72\%.} \\
\texttt{5. For students who are encouraged, the probability of high exam score is 52\%.} \\
\\ 
\texttt{We are asked to determine whether high exam score is more likely than low exam score overall. To do this, we need to consider the potential outcomes under different scenarios.} \\
\\ 
\texttt{Let's define two potential outcomes for each student:} \\
\\ 
\texttt{{-} High exam score (H)} \\
\texttt{{-} Low exam score (L)} \\
\\ 
\texttt{We can also define two counterfactuals:} \\
\\ 
\texttt{{-} High encouragement (E+)} \\
\texttt{{-} No encouragement (E{-})} \\
\\ 
\texttt{Using the given information, we can estimate the probabilities of these potential outcomes:} \\
\\ 
\texttt{{-} P(E+) = 0.91 (probability of high encouragement)} \\
\texttt{{-} P(E{-}) = 0.09 (probability of no encouragement)} \\
\texttt{{-} P(H|E{-}) = 0.72 (probability of high exam score given no encouragement)} \\
\texttt{{-} P(H|E+) = 0.52 (probability of high exam score given high encouragement)} \\
\\ 
\texttt{Now, we need to consider the potential outcomes under the counterfactual scenarios. Specifically, we want to know the probability of high exam score (H) under the scenario where a student is not encouraged (E{-}). We can use the law of total probability to estimate this:} \\
\\ 
\texttt{P(H|E{-}) = P(H|E{-}) \textbackslash{}* P(E{-}) + P(H|E+) \textbackslash{}* P(E+)} \\
\texttt{= 0.72 \textbackslash{}* 0.09 + 0.52 \textbackslash{}* 0.91} \\
\texttt{= 0.0648 + 0.4722} \\
\texttt{= 0.537} \\
\\ 
\texttt{Similarly, we can estimate the probability of low exam score (L) under the scenario where a student is not encouraged (E{-}):} \\
\\ 
\texttt{P(L|E{-}) = 1 {-} P(H|E{-}) = 1 {-} 0.537 = 0.463} \\
\\ 
\texttt{Now, we can estimate the overall probability of high exam score (H) by considering the probability of high encouragement (E+) and the counterfactual scenario where a student is not encouraged (E{-}):} \\
\\ 
\texttt{P(H) = P(E+) \textbackslash{}* P(H|E+) + P(E{-}) \textbackslash{}* P(H|E{-})} \\
\texttt{= 0.91 \textbackslash{}* 0.52 + 0.09 \textbackslash{}* 0.537} \\
\texttt{= 0.4752 + 0.0483} \\
\texttt{= 0.5235} \\
\\ 
\texttt{Similarly, we can estimate the overall probability of low exam score (L):} \\
\\ 
\texttt{P(L) = 1 {-} P(H) = 1 {-} 0.5235 = 0.4765} \\
\\ 
\texttt{Since P(L) > P(H), we conclude that low exam score is more likely than high exam score overall.} \\
            }
        }
    }
    & & \\ \\

    \theutterance \stepcounter{utterance}  
    & & & \multicolumn{2}{p{0.3\linewidth}}{
        \cellcolor[rgb]{0.9,0.9,0.9}{
            \makecell[{{p{\linewidth}}}]{
                \texttt{\tiny{[GM$|$GM]}}
                \texttt{No} \\
            }
        }
    }
    & & \\ \\

    \theutterance \stepcounter{utterance}  
    & & & \multicolumn{2}{p{0.3\linewidth}}{
        \cellcolor[rgb]{0.9,0.9,0.9}{
            \makecell[{{p{\linewidth}}}]{
                \texttt{\tiny{[GM$|$GM]}}
                \texttt{yes} \\
            }
        }
    }
    & & \\ \\

    \theutterance \stepcounter{utterance}  
    & & & \multicolumn{2}{p{0.3\linewidth}}{
        \cellcolor[rgb]{0.9,0.9,0.9}{
            \makecell[{{p{\linewidth}}}]{
                \texttt{\tiny{[GM$|$GM]}}
                \texttt{game\_result = LOSE} \\
            }
        }
    }
    & & \\ \\

\end{supertabular}
}

\end{document}
